\documentclass[11pt]{article}
\usepackage[margin=1in]{geometry}
\usepackage{amsmath,amssymb,booktabs}
\usepackage[hidelinks]{hyperref}
\title{Covariant host stress, density dependence, and finite-stock tails}
\author{Zimmerman programme: CPB-MF-2026-10-16}
\date{16 October 2026}
\begin{document}
\maketitle
\section{Result and scope}
The missing frame stress is physically consequential. A band sea whose scales
are fixed by the auxiliary flux can be a vacuum potential in a chosen
covariant completion. A sea whose scales follow material compression instead
has radiation-like stress; its negative energy cannot by itself support a
healthy perfect fluid. An explicit positive host repairs the elementary
longitudinal signs, but supplies new equation-of-state data. Neither choice
identifies the host with the original MeV particles.

A density compensation can preserve the leading cubic response. It transfers
density dependence into the quadratic critical subtraction and, in the
regulated local toy, worsens the infinite-tail energy. A finite material
inventory alone therefore does not complete the halo. These are conditional
calculations on specified effective actions, not a full covariant fermion,
metric, formation or quantum completion. Units are $\hbar=c=1$, metric
signature $(-+++)$, fixed positive flux $Q$ of mass dimension two.

\section{A covariant barotropic host}
Introduce three scalar material labels $X^I$ and define
\begin{align}
 B^{IJ}&=g^{\mu\nu}\partial_\mu X^I\partial_\nu X^J,&
 b&=\sqrt{\det B}>0,\\
 J^\mu&=\frac{\epsilon^{\mu\nu\rho\sigma}\epsilon_{IJK}}
 {3!\sqrt{-g}}\partial_\nu X^I\partial_\rho X^J\partial_\sigma X^K,
 &u^\mu&=J^\mu/b,\qquad u^\mu u_\mu=-1.
\end{align}
Orientation selects a future-directed current. Antisymmetry gives
$\nabla_\mu J^\mu=0$ identically. Choose
\begin{equation}
 S_h=-\int d^4x\sqrt{-g}\,\rho(b,Q).
\end{equation}
The standard nondissipative fluid dictionary is derived in
Ref.~\cite{MFfluid}, especially its equations (11)--(13), (25)--(34).
Its charged-fluid thermodynamic convention identifies $b$ with entropy
density. Here a separately normalized barotropic material current is used;
a relation to number requires fixed-composition, isentropic microscopic
matching. It is not the signed fermion density $q_F$, the carrier charge
$Q_c$, or an established cold-particle number.

Varying the metric at fixed $Q$ gives
\begin{equation}
 T_{\mu\nu}^{h}=(\rho+p)u_\mu u_\nu+p g_{\mu\nu},\qquad
 p=b\rho_b-\rho,\qquad \rho+p=b\rho_b.
 \label{mf:stress}
\end{equation}
The pressure also follows from $-\partial_V[V\rho(N/V,Q)]_N$.
For homogeneous fixed-background perturbations, up to a boundary term,
\begin{equation}
 \mathcal L_h^{(2)}=\frac{\rho+p}{2}
 [\dot{\boldsymbol\pi}^{,2}-c_s^2(\boldsymbol\nabla\cdot
 \boldsymbol\pi)^2],\qquad c_s^2=\frac{b\rho_{bb}}{\rho_b}.
 \label{mf:phonon}
\end{equation}
A direct one-dimensional check uses
$b=b_0\sqrt{(1+\epsilon\partial_x\pi)^2-\epsilon^2\dot\pi^2}$.
Thus $\rho+p>0$ and $\rho_{bb}>0$ are sufficient for positive temporal
and longitudinal gradient terms. Transverse gradients vanish at this
perfect-fluid order. This is not a transverse quantum, full fermion,
metric or ultraviolet stability certificate. At $b=0$ the material frame
is undefined; a vanishing dilute enthalpy does not establish a uniformly
controlled vacuum phase.

\section{The filled sea has two different stress interpretations}
Let $\mu=\sqrt Q$ and use the previous example's positive constants
$A_*=1.8006326323142123$, $c_\Lambda=c_h=1/2$. Define
\begin{equation}
 y=b/Q^{3/2},\qquad d=A_*c_\Lambda^3/3
 =0.07502635967975885.
\end{equation}
If $A=A_*\mu$ and $\Lambda_F=c_\Lambda\mu$ are independent of $b$,
the sea term $\epsilon_{\rm sea}=-A\Lambda_F^3/3=-dQ^2$ is a
$b$-independent potential in this chosen completion. Equation
\eqref{mf:stress} gives $p=-\epsilon_{\rm sea}$. Its vacuum-form
metric stress may enter the previous sequestering projection, while its
$Q$ derivative must remain in the electric constraint. This observation
does not prove that an actual occupied band and its frame have only that
stress.

If all momentum scales instead compress as $b^{1/3}$ with fixed velocity,
then $A=A_*\mu y^{1/3}$ and $\Lambda_F=c_\Lambda\mu y^{1/3}$ imply
\begin{equation}
 \epsilon_{\rm sea}=-d b^{4/3},\qquad
 p_{\rm sea}=\epsilon_{\rm sea}/3,\qquad
 T^{\mu}_{\ \mu,\rm sea}=0,\qquad
 \rho_{\rm sea}+p_{\rm sea}=-\tfrac43d b^{4/3}<0.
\end{equation}
This term survives subtraction of a vacuum-form constant: its trace is
zero and its stress is nonzero. Adding a constant vacuum potential cannot
repair its kinetic sign. Adding a rest term $M(Q)b$ can make the enthalpy
positive at low density, but $\rho_{bb}=-4d b^{-2/3}/9<0$ then gives a
longitudinal instability. A positive finite $b^2$ term fails to dominate
that curvature as $b\to0$.

For example, writing $\rho/Q^2=c_0+c_1 y+c_2 y^2/2-d y^{4/3}$,
its curvature changes sign at
\begin{equation}
 y_{\rm crit}=(4d/(9c_2))^{3/2}.
\end{equation}
For $c_2=1$ this is $0.00608901487444$. Positivity above this threshold
still needs the enthalpy test. It is not a theorem excluding other hosts.

An explicit sufficient repair is to supply a positive host $y^{4/3}$
coefficient $2d$, leaving net $\lambda=d>0$ after the sea:
\begin{equation}
 \rho=Q^2 f(y),\qquad
 f=c_0+c_1 y+\lambda y^{4/3}+c_2y^2/2,
 \qquad c_1,\lambda,c_2>0.
 \label{mf:healthy}
\end{equation}
For every $y>0$, $f'>0$, $f''>0$, and
$f'-yf''=c_1+(8\lambda/9)y^{1/3}>0$, proving
$\rho+p>0$ and $0<c_s^2<1$. Its trace is
\begin{equation}
 -\rho+3p=Q^2[-4c_0-c_1 y+c_2 y^2].
\end{equation}
The extra coefficients are inputs requiring particle/state matching.
The illustrative $c_0=0,c_1=c_2=1,\lambda=d$ was checked at 61 densities
from $10^{-9}$ to $10^3$ in addition to these exact identities.

\section{Cosmological roles and spectral compensation}
In a homogeneous comoving phase, current conservation fixes
$b\propto a^{-3}$ at constant $Q$. Consequently the four terms of
\eqref{mf:healthy} redshift as $a^0,a^{-3},a^{-4},a^{-6}$ respectively,
and obey $\dot\rho+3H(\rho+p)=0$. Only the constant has vacuum stress.
No vacuum-to-material production follows from this identity. Heating,
entropy production and non-equilibrium particle creation require additional
kinetics beyond this nondissipative action.

The compensated band excess has the exact local regulated expression
\begin{align}
 V_F&=-\frac A3[(\Lambda_F^2+\Delta_f^2)^{3/2}
             -\Lambda_F^3-\Delta_f^3]
             +\frac{A\Lambda_F\Delta_f^2}{2},\\
 \Delta_f&=hp,\quad p=|\boldsymbol P|,\qquad
 V_F=\frac\beta3p^3+O(p^4),\quad \beta=Ah^3.
 \label{mf:det}
\end{align}
Its quadratic critical counterterm is $K_{\rm ct}p^2$ with
$K_{\rm ct}=A\Lambda_Fh^2/2$. Three explicitly chosen assignments are
\begin{center}
\begin{tabular}{lccccc}
\toprule
Choice & $A/\mu$ & $\Lambda_F/\mu$ & $h\mu$ & $\beta$ scaling
 & $K_{\rm ct}$ scaling\\
\midrule
Flux & $A_*$ & $c_\Lambda$ & $c_h$ & $Q^{-1}$ & $1$\\
Material & $A_*y^{1/3}$ & $c_\Lambda y^{1/3}$ & $c_hy^{-1/3}$
 & $Q^{-1}y^{-2/3}$ & $1$\\
Compensated cubic & $A_*y^{1/3}$ & $c_\Lambda y^{1/3}$ & $c_hy^{-1/9}$
 & $Q^{-1}$ & $y^{4/9}$\\
\bottomrule
\end{tabular}
\end{center}
For the previous fixed-coupling annular force dictionary, $a_0\propto
\beta^{-1}$. Thus the material assignment gives $a_0\propto a^{-2}$
in a homogeneous expanding phase, rather than a constant vacuum-linked
acceleration. A bound host can have another density history; this is a
conditional mismatch, not a general exclusion of the theory.

At fixed velocity, $A,\Lambda_F\propto b^{1/3}$ and $h\propto b^\gamma$
give exponents $1/3+3\gamma$ for $\beta$ and $2/3+2\gamma$ for
$K_{\rm ct}$. Their zeros are incompatible. Changing other spectral data
can evade this restricted result. The compensated cubic assignment selects
$\gamma=-1/9$ and requires the critical subtraction itself to vary with
material density. In particular
\begin{equation}
 b\partial_b(K_{\rm ct}p^2)=\tfrac49 K_{\rm ct}p^2.
\end{equation}
Omitting this contribution after matching the force changes the host
Euler/chemical equation. The total finite-cutoff excess has a small
nonzero rearrangement even when its leading cubic coefficient is constant.
For $Q=1,p=10^{-4},y=1$, it is
$b\partial_bV_F/V_F=1.66672916\times10^{-5}$, whereas the separate
counterterm contribution is $5.00175731\times10^{-10}$ and
$V_F=7.50235462\times10^{-14}$. The cancellation must be varied as a
whole. At nonzero $\boldsymbol P$, this derivative is a material
rearrangement at fixed $p$, not the full isotropic pressure: the norm,
frame and polarization also vary with the metric. Only the $P=0$ host
stress has been completed here.

\section{Finite inventory does not by itself cure the tail}
Prescribe, for diagnosis only, a spherical host
\begin{equation}
 y(r)=(1+r^2)^{-2},\quad \int d^3x\,y=\pi^2,\qquad
 D(r)=r/(1+r^2)^{3/2},\qquad C=Q=1.
\end{equation}
This is a finite-count profile with a chosen unit Plummer source, not a
solution of the coupled host or gravitational equations. The exact local
force equation from \eqref{mf:det} is
\begin{equation}
 D=\frac{\partial V_F}{\partial p}
 =\frac{2Ah\Lambda_F\Delta_f^2}
 {\Delta_f+\Lambda_F+\sqrt{\Delta_f^2+\Lambda_F^2}}.
 \label{mf:tail}
\end{equation}
It was solved on 161 radii $1\le r\le10^8$ per choice, with relative
force tolerance $10^{-9}$ and 60-digit energy evaluation. Fits over the
last 20 points give the asymptotic regulated-toy exponents
\begin{center}
\begin{tabular}{lccc}
\toprule
Choice & $p(r)$ & $dE_F/d\ln r=4\pi r^3V_F$ & low-gap tail retained?\\
\midrule
Flux & $r^{-1}$ & $r^0$ & yes\\
Material & $r^{-2}$ & $r^{-1}$ & no\\
Compensated cubic & $r^{-2/9}$ & $r^{7/9}$ & no\\
\bottomrule
\end{tabular}
\end{center}
The fitted exponents agree to better than $10^{-6}$ (gate $0.01$).
The low-gap flag means $\Delta_f/\Lambda_F<0.1$ at the last sample;
false flags are retained. The high-gap extensions describe this regulated
local expression only, not a controlled microscopic band theory.

For the flux choice the persistent cubic $p\propto r^{-1}$ gives a
logarithmically divergent excess energy despite finite host count. For
the material choice the high-gap expression becomes
$V_F\sim A\Lambda_Fh^2p^2/2$ with constant $A\Lambda_Fh^2$;
the toy returns to an inverse-square response and its tail energy converges.
Its singular coupling and vanishing cutoff as $b\to0$ require a new phase
completion before claiming a physical finite halo.

For the compensated choice on a general finite-mass power tail
$y\sim r^{-k}$, $k>3$, a putative cubic branch has
$\Delta_f/\Lambda_F\sim r^{-1+4k/9}$ and necessarily leaves the
small-gap regime. Continuing the regulated toy gives
\begin{equation}
 A\Lambda_Fh^2\sim r^{-4k/9},\quad
 p\sim r^{-2+4k/9},\quad
 \frac{dE_F}{d\ln r}\sim r^{-1+4k/9},
\end{equation}
which diverges. This is a conditional asymptotic obstruction for the
specified positive local band energy; it does not include a solved
interface, host relaxation, nonlocal band gradients or gravitational
binding energy and is not a universal no-go theorem.

\section{Research consequence and provenance}
A useful next route is a microscopically selected finite-density phase
with an actual interface and a controlled outside phase. A density plateau
could hold coefficients fixed over an annulus, but its density, charge,
inventory, stresses, interface energy and source-dependent formation must
all follow from one action. Merely prescribing a cutoff or profile does
not supply these conditions. The material-frame action here is a concrete
piece of that task, not the missing full fermion action.

The earlier scalar carrier, neutral production network and charge-even
portal remain distinct candidates requiring identification with this
host. Their abundance, mediator populations, narrow radiative thresholds
and transport obligations remain open. The auxiliary residual flux/history
condition, vacuum magnitude, fitted $\kappa$, finite-stock galaxy dynamics,
relativistic lensing and cosmological data tests also remain open. No
full physical closure or globally novel mechanism is claimed.

Evidence: \texttt{material\_frame\_audit\_run\_001} contains 156 controls;
\texttt{material\_tail\_probe\_run\_002} contains 972 controls, including
483 force roots and positive-energy checks. Both manifests passed freshness
validation. Tail run 001 is retained as failed: a NumPy boolean prevented
JSON serialization after calculation. Run 002 converts that output flag
to a built-in boolean; no physical gates or tolerances changed. Exact
identities, finite numerical ranges and out-of-regime flags are separated
above. This is self-review, without an independent referee or global
novelty search. Frozen theory criteria were not changed.

\begin{thebibliography}{9}
\bibitem{MFfluid} S. Dubovsky, L. Hui, A. Nicolis and D. T. Son,
\emph{Effective field theory for hydrodynamics: thermodynamics, and the
derivative expansion}, arXiv:1107.0731v1 (4 July 2011),
\url{https://arxiv.org/html/1107.0731v1}. Version identified by arXiv's
header; the generated internal typesetting date is not used as the
publication date. The primary fluid dictionary is established literature,
not an originality claim for this programme.
\end{thebibliography}
\end{document}
